\documentclass[10pt,a4paper]{amsart}
\usepackage[margin=19mm]{geometry}
\usepackage{amsmath,amssymb,mathtools}
\usepackage[scr=rsfs,cal=euler]{mathalfa}
\usepackage{xfrac}
\usepackage[unicode,hidelinks,bookmarks=false]{hyperref}
\newcommand{\ie}{i.e.\ }
\newcommand{\cf}{cf.\ }
\newcommand{\resp}{respectively\ }
\newcommand{\Subsection}{\subsection}
\providecommand{\nobreakdash}{\nobreak\hbox{-}}
\setlength{\emergencystretch}{2em}
\multlinegap0pt
\usepackage{bm}
\usepackage{bbm}
\usepackage{dsfont}
\usepackage{xspace}
\usepackage{cancel}
\usepackage{mathtools}
\usepackage{amsthm}
\makeatletter
\@ifundefined{theorem}{\newtheorem{theorem}{Theorem}}{}
\makeatother
\newcommand{\Cc}{\mathcal{CC}^{\mathrm{onf}}}
\title[On the Configuration Space of planar closed Kinematic Chains]{On the Configuration Space of planar closed Kinematic Chains}
\author{Gerhard Zangerl}
\date{Source: arXiv:1902.10824v3 (CC BY 4.0)}
\begin{document}
\maketitle

\noindent\emph{Source and licence.} On the Configuration Space of planar closed Kinematic Chains. Version arXiv:1902.10824v3 (CC BY 4.0), distributed under the Creative Commons Attribution 4.0 International licence, as recorded in the arXiv licence metadata for this exact version and shown on the abstract page https://arxiv.org/abs/1902.10824v3. Full rights notice: licenses/arxiv-1902.10824-CC-BY-4.0.txt.\par\smallskip

\noindent\textit{Editorial scope.} Counted unit: the Theorem whose source label is \texttt{thm:sampling} in the article (source0.tex lines 327--338, source label \texttt{thm:sampling}), delivered together with the article's own complete original proof of it (source0.tex lines 340--370, one \texttt{proof} environment of 31 lines). No mathematics is written, repaired or re-derived: statement and proof are the article's own text line for line. Required context retained uncounted: eq:addsincos (lines 235--237); eq:SQshort (lines 311--313). Every definition, display and label of those lines is retained unchanged. Omitted from this package and not redistributed: 1--234, 238--310, 314--326, 339--339, 371--635. Nothing inside a retained excerpt is omitted or altered: every retained body is delivered exactly as the article's lines are written, so no omission receipt of any kind accompanies this package. Citations: the retained excerpts carry no citation command at all, so no bibliography entry of the article is transcribed and no reference is redistributed. Reference adaptation: none was needed and none was made; every reference inside the delivered unit resolves inside it, so no unresolved reference or citation is produced. Printed slips kept literal: no printed word, symbol, hypothesis or number of the counted statement or of its proof was altered. Layout and class: the article's own class is \texttt{article} with packages including babel, fontenc, inputenc, arev, geometry, amsmath, amssymb, amsthm, graphicx, xcolor, adjustbox, tikz, pgfplots; the wrapper is the lane's pinned \texttt{amsart} editorial wrapper at 10pt on A4, which restates the theorem-like environments the retained text needs and adds the article's own macro definitions that the retained text needs (\texttt{\textbackslash Cc}), with extra wrapper packages (bm, bbm, dsfont, xspace, cancel, mathtools). The delivered numbering is the unit's own. Assets: no retained span contains a figure environment, a graphics-inclusion command, a PDF-page inclusion, a table or a TikZ picture; no image file is copied into this package, the package declares no asset dependency and no image file is redistributed with it. Licence: version arXiv:1902.10824v3 of this article is distributed under the Creative Commons Attribution 4.0 International licence, as recorded in the arXiv licence metadata for this exact version and shown on the abstract page https://arxiv.org/abs/1902.10824v3. A full rights notice is delivered at licenses/arxiv-1902.10824-CC-BY-4.0.txt. Characters: every retained excerpt is pure ASCII, so the delivered engine, XeLaTeX with its default Latin Modern text font, reports no missing character.
\begin{align}\label{eq:addsincos}
a \sin\left( x \right) + b \cos\left( x\right) = c \sin\left(x + \varphi(a,b)  \right), 
\end{align}

\begin{align}\label{eq:SQshort}
2 a_{n-1} \sin\left( \beta_{n-1} \right) Y_{n-2}\left( \beta^{n-2} \right)  &+ 2a_{n-1} \cos\left( \beta_{n-1} \right) X_{n-2}\left( \beta^{n-2} \right) +  L_{n-2}^2 =  L_{n-1}^2 - a_{n-1}^2.  
\end{align}

\begin{theorem}[Computation of circular configurations]\label{thm:sampling} Every circular configuration
\(\beta^{n-1}\in  \Cc(a^n)\)
has a diagonal vector contained in
\(\mathcal{DS}(a^n)\).
Conversely, for every
\(L^{n-3}\in\mathcal{DS}(a^n)\),
there exists a circular configuration having these diagonal lengths. Moreover, the the angle $\beta_k$ is related to the angles $\beta^{k-1}$ by the formula 
\begin{align}\label{eq:diagangle}
2 a_k L_{k-1}  \sin\left( \beta_k + \Phi_{k-1}\left( \beta^{k-1}\right)  \right) =  L_k^2  - a_k^2 - L_{k-1}^2,   
\end{align}
for $1 \leq k \leq n-1$, whenever $\Phi_{k-1}$ is defined. If not, the relation between $\beta_k$ and $\beta^{k-1}$ is given by \eqref{eq:SQshort}, where $n-1$ and $n-2$ are replaced by $k$ and $k-1$. Note that for $k=1$ we set $L_0 := 0$ and $\Phi_0 :=0$ so that $\beta_1$ is an arbitrary angle. 
\end{theorem}

\begin{proof} Assume $\beta^{n-1}$ is a circular configuration and thus solves equation \eqref{eq:SQshort}. We will manipulate equation \eqref{eq:SQshort} to show that the vector of variable diagonal lengths $L^{n-2}$ is indeed an element of $\mathcal{DS}({a^n})$. We start by showing  that the diagonal $L_{n-2}$ is contained  in $[|L_{n-1}-a_{n-1}|, L_{n-1}+a_{n-1}] \cap [ \mathrm{    R}_{n-2}^{\min} , \mathrm{    R}_{n-2}^{\max}]$. For the remaining diagonals the analogous statement for $L_{n-k-1}$ then follows inductively by applying the same arguments.\\

Although it is very tempting to apply addition formula \eqref{eq:addsincos} to this equation we have to  deal with special cases before we can do so.\\

$X_{n-2} \neq 0$ but $Y_{n-2} = 0$. In this case, according to \eqref{eq:SQshort} we obtain the equation  \begin{align*}
& 2 a_{n-1} \cos(\beta_{n-1} )  X_{n-2} + L_{n-2}^2  =  L_{n-1}^2 - a_{n-1}^2.  
\end{align*}
solving for $\beta_{n-1}$ yields  
\begin{align*}
& \cos(\beta_{n-1} )     =  \frac{ L_{n-1}^2 - L_{n-2}^2- a_{n-1}^2 }{ 2 a_{n-1}  X_{n-2}  }.  
\end{align*}
The Latter equation can be solved for $\beta_{n-1}$ iff the square of its right side is lower or equal to one. Since in the considered case $ X_{n-2}^2 = L_{n-2}^2$,  this leads to the inequality   
\begin{align}\label{ieq:diagonals}
\left(   L_{n-1}^2 - L_{n-2}^2- a_{n-1}^2 \right)^2 \leq  4 a_{n-1}^2 L_{n-2}^2, 
\end{align}
that is satisfied for $L_{n-2} \in [ |L_{n-1} -a_{n-1}|,  L_{n-1} +a_{n-1} ]$. Additionally, since $L_{n-2}$ is a diagonal length also $L_{n-2} \in [\mathrm{R}_{n-2}^{\min}, \mathrm{    R}_{n-2}^{\max}]$ has to be satisfied. The case when $X_{n-2} = 0$ and $Y_{n-2} \neq 0$ is analogously.\\       

The case $ L_{n-2} = X_{n-2} = Y_{n-2} = 0$ can only  occur for $a_n = L_{n-1} = a_{n-1}$. Here $\beta_{n-1}$ is an arbitrary value, and $\mathrm{R}_{n-2}^{\min} = 0$ since $L_{n-2} = 0 $  is a diagonal of the circular configuration. Also $0 = L_{n-2} \in [ |L_{n-1} -a_{n-1}|,  L_{n-1} +a_{n-1} ] = [0, 2 a _n ]$ is satisfied.\\ 

If $X_{n-2} \neq 0 $ and  $Y_{n-2} \neq 0 $ we can apply the addition theorem to formula \eqref{eq:SQshort} and obtain   
\begin{align*}
& 2 a_{n-1} \sqrt{  X_{n-2}^2 + Y_{n-2}^2 } \, \sin\left( \beta_{n-1} + \Phi_{n-2}\left(\beta^{n-2} \right)  \right) + L_{n-2}^2  = L_{n-1}^2 - a_{n-1}^2,   
\end{align*}
which can be rewritten as 
\begin{align*}
& 2 a_{n-1} L_{n-2} \sin\left( \beta_{n-1} + \Phi_{n-2}\left(\beta^{n-2} \right)  \right) + L_{n-2}^2  = L_{n-1}^2 - a_{n-1}^2.   
\end{align*}
Repeating the arguments outlined above We conclude  that this equation can only have a solution if \eqref{ieq:diagonals} is satisfied and thus also in this case $L_{n-2}  \in  [ |L_{n-1} -a_{n-1}|,  L_{n-1} +a_{n-1} ]$ and $L_{n-2} \in [\mathrm{R}_{n-2}^{\min}, \mathrm{    R}_{n-2}^{\max}]$. Thus, repeating the arguments it follows that $L^{n-3} \in \mathcal{DS}(a^n)$.\\
       
Conversely, if $L^{n-3} \in \mathcal{DS}(a^n)$ then  equation \eqref{eq:diagangle} can be solved for $\beta_k$ if $\Phi_{k-1}$ is defined. If not, equation \eqref{eq:SQshort} can be used alternatively, with $n-1$, $n-2$ replaced by $k$ and $k-1$. Repeating this procedure gives a circular configuration $\beta^{n-1}$. 
\end{proof}

\end{document}
